\section{Experimental Setup}
\label{section:experiments}

This section describes our benchmark dataset, training hyperparameters, the ablation study design, and the hardware profiling methodology used to evaluate Pyramid-ResMLP.

\subsection{Benchmark dataset: Tiny-ImageNet-200}
\label{subsection:dataset}
We evaluate all model variants on the \textbf{Tiny-ImageNet-200} benchmark~\cite{le2015tiny}. Tiny-ImageNet is a challenging subset of the full ImageNet dataset, consisting of:
\begin{itemize}
    \item 200 distinct object classes;
    \item 100,000 training images (500 images per class);
    \item 10,000 validation images (50 images per class);
    \item Fixed native resolution of $64 \times 64$ pixels with 3 color channels.
\end{itemize}
Because the image resolution is small ($64 \times 64$) yet contains 200 diverse categories, models cannot rely on memorizing high-frequency spatial patterns. Instead, they must extract meaningful structural features, making it an ideal testbed for evaluating the inductive bias and parameter efficiency of vision MLPs.

\subsection{Training configurations and regularization}
\label{subsection:training_setup}
All experiments are executed in a distributed training setup on dual NVIDIA Tesla T4 GPUs with an effective batch size of 128 for 60 epochs. We use the AdamW optimizer~\cite{loshchilov2017adamw} with decoupled weight decay.

The complete set of training hyperparameters is detailed in Table~\ref{tab:hyperparams}.

\begin{table}[htbp]
\centering
\small
\adjustbox{max width=\linewidth}{%
\begin{tabular}{ll}
\toprule
\textbf{Hyperparameter} & \textbf{Value} \\
\midrule
Input Resolution & $64 \times 64 \times 3$ \\
Epochs & 60 \\
Batch Size & 128 (Distributed across 2 $\times$ Tesla T4) \\
Optimizer & AdamW ($\beta_1 = 0.9, \beta_2 = 0.999$) \\
Base Learning Rate ($\eta_0$) & $1.0 \times 10^{-3}$ \\
Min Learning Rate ($\eta_{\text{min}}$) & $1.0 \times 10^{-6}$ \\
Learning Rate Schedule & Cosine Annealing~\cite{loshchilov2016sgdr} with 5 warmup epochs \\
Weight Decay & $0.05$ \\
Loss Function & Cross-Entropy with Label Smoothing ($\epsilon = 0.1$) \\
Data Augmentation & RandAugment ($N=2, M=9$)~\cite{cubuk2020randaugment}, Random Horizontal Flip ($p=0.5$) \\
Mixup ($\alpha$) & $0.8$~\cite{zhang2017mixup} \\
CutMix ($\alpha$) & $1.0$~\cite{yun2019cutmix} \\
DropPath (Stochastic Depth) & $0.1$~\cite{huang2016deep} \\
Gradient Clipping & Maximum norm $1.0$ \\
\bottomrule
\end{tabular}%
}
\caption{Training hyperparameters and regularization pipeline on Tiny-ImageNet-200.}
\label{tab:hyperparams}
\end{table}

\subsection{Ablation study matrix}
\label{subsection:ablation_matrix}
To isolate the exact contribution of each architectural innovation, we conduct a controlled four-way ablation study where all models share identical training data, optimizer settings, learning rate schedules, and seeds:
\begin{enumerate}
    \item \textbf{Ablation 1 (pure linear baseline):} Pyramid-ResMLP Pure with an overlapping stem and pure linear token-mixing (\texttt{reparam\_conv\_k: []}). This isolates the performance of hierarchical downsampling without any convolutional inductive bias.
    \item \textbf{Ablation 2 (DW $3 \times 3$ only):} Adds only the $3 \times 3$ depthwise convolution branch (\texttt{reparam\_conv\_k: [3]}) to the overlapping stem. This isolates the specific benefit of local 2D spatial context.
    \item \textbf{Ablation 3 (full RepTokenMix):} Integrates both $1 \times 1$ and $3 \times 3$ depthwise convolution branches (\texttt{reparam\_conv\_k: [1, 3]}) with an overlapping stem. This is our complete proposed design.
    \item \textbf{Ablation 4 (non-overlapping stem):} Replaces the overlapping stem with a non-overlapping $4 \times 4$ convolutional projection (kernel 4, stride 4, padding 0), while keeping Full RepTokenMix (\texttt{reparam\_conv\_k: [1, 3]}). This isolates the specific impact of stem patch continuity versus disjoint patch slicing.
\end{enumerate}

\subsection{Hardware inference profiling protocol}
\label{subsection:profiling_protocol}
To evaluate execution latency, throughput, and memory consumption under realistic production conditions, we profile the model on an NVIDIA Tesla T4 GPU (Turing architecture, 16~GB VRAM) under FP32 precision.

The profiling protocol is structured as follows:
\begin{itemize}
    \item \textbf{Batch sizes tested:} $B \in \{1, 16, 32, 64, 128\}$.
    \item \textbf{Warmup phase:} 50 iterations of forward passes to stabilize GPU operating clocks, warm up CUDA caches, and eliminate driver initialization overhead.
    \item \textbf{Measurement phase:} 200 consecutive forward iterations timed using asynchronous CUDA events (\texttt{torch.cuda.Event}) with device synchronization (\texttt{torch.cuda.synchronize()}).
    \item \textbf{Memory measurement:} Peak allocated GPU VRAM is measured using PyTorch's \texttt{max\_memory\_allocated()} API, with the CUDA cache flushed prior to each profiling run.
    \item \textbf{Numerical equivalence check:} We calculate the Chebyshev ($\ell_\infty$) error norm between training-mode multi-branch outputs and deployed fused-matrix outputs across all tokens:
    \begin{equation}
        \Delta_{\text{max}} = \max_{b, n, c} \big| Y_{\text{train}}[b, n, c] - Y_{\text{deploy}}[b, n, c] \big|
    \end{equation}
    An equivalence check is considered passed if $\Delta_{\text{max}} < 1.0 \times 10^{-6}$ (within standard 32-bit floating-point precision limits).
\end{itemize}
